\section{Conclusion}
\label{sec:conclusion}

We formulated garment pattern layout as a joint optimisation over seam
positions, fabric efficiency, and texture alignment, with an explicit
fit-preservation guard, and demonstrated its solvability with a
mixed-integer evolution strategy.
By optimising seam positions directly on the 3D body mesh, all pattern
modifications are grounded in body geometry, making fit deviation an
explicit and measurable quantity.
A global phase alignment stage and discrete texture phase offsets per patch
exploit the full periodic symmetry of the fabric, modelled across eight
wallpaper groups including those with glide-reflection symmetry.
Systematic ablation decomposes the GA's advantage into three individually
attributable components: $\boldsymbol{\delta}$ search provides the
largest single improvement (12\% over the best $\boldsymbol{\kappa}$-only
baseline), evolutionary operators exploit the sparse
$\boldsymbol{\delta}$ landscape more effectively than random search
(${\sim}$3\%), and $\boldsymbol{\kappa}$ mutation contributes a further
${\sim}$6\% by reaching global phase reconfigurations inaccessible to
Stage~2 alone.
Experiments on three garment types ($M \in \{2, 4, 11\}$) with up to
2{,}200 pattern pieces demonstrate consistent improvement in texture
alignment over all baselines including random search at matched
evaluation budget.

\paragraph{Limitations and future work.}
The nesting engine uses a bottom-left heuristic; substituting an
NFP-based solver~\cite{MartinezSykoraACRO22} could improve fabric
utilisation without affecting the other contributions.
Approximately 55--60\,\% of candidate $\boldsymbol{\delta}$ configurations
fail during geometry processing (geodesic or LSCM singularities),
receiving death-penalty fitness; more robust geometry processing or
adaptive step-size control would improve effective population utilisation.
The discrete Voronoi structure of the $\boldsymbol{\delta}$ landscape
(Section~\ref{sec:ea}) causes many neutral mutations; adaptive step-size
control or landscape-aware operators could accelerate convergence.
Finally, extending the framework to multiple body shapes and sizes within
a single fabric roll is a natural next step toward industrial deployment.
